% Generated draft: requires mathematical and source review.
\documentclass[11pt]{article}
\usepackage[margin=25mm]{geometry}
\usepackage{fontspec}
\setmainfont{Latin Modern Roman}
\usepackage{amsmath,amssymb,mathtools,unicode-math}
\setmathfont{Latin Modern Math}
\usepackage{enumitem,needspace,xurl,hyperref}
\setlength{\parindent}{0pt}
\setlength{\parskip}{5pt}
\setlength{\emergencystretch}{3em}
\newcommand{\N}{\mathbb{N}}
\newcommand{\Z}{\mathbb{Z}}
\newcommand{\Q}{\mathbb{Q}}
\newcommand{\R}{\mathbb{R}}
\newcommand{\C}{\mathbb{C}}
\newcommand{\F}{\mathbb{F}}
\newcommand{\A}{\mathbb{A}}
\newcommand{\PP}{\mathbb{P}}
\newcommand{\E}{\mathbb{E}}
\newcommand{\eps}{\varepsilon}
\newcommand{\dd}{\,\mathrm{d}}
\begin{document}
\section*{UnsolvedMath ID: 20001975}
\textbf{Source problem code:} catalogue entry
\subsection{Definitions and mathematical statement}
Problem 10.1, attributed to R. Goldin, S. Tolman, and J. Weitsman, asks for the conditions under which Kirwan surjectivity for symplectic reduction holds over \(\Z\), and for the most general such conditions. It is a classification request; the original excerpt does not fix compactness, group connectedness, level regularity, freeness, or whether a nonfree target is Borel, orbifold, or coarse-space cohomology.

To state the supplied progress in a standard regular-level setting, let \((M,\omega)\) be a compact Hamiltonian manifold for a compact connected Lie group \(G\), with equivariant moment map \(\mu:M\to\mathfrak g^*\). Assume \(0\) is a regular value, set \(Z=\mu^{-1}(0)\), and let \(i:Z\hookrightarrow M\). The Borel Kirwan restriction with coefficients in a commutative ring \(R\) is
\[
 \kappa_R^{\mathrm B}=i^*:H_G^*(M;R)\longrightarrow H_G^*(Z;R),\qquad
 H_G^*(Y;R)=H^*(EG\times_GY;R).
\]
If \(G\) acts freely on \(Z\), the usual comparison identifies the target with \(H^*(Z/G;R)\). The aid attributes to Pomerleano--Teleman additive split surjectivity over \(\Z\) in the free regular-level case. A splitting is additive, not an asserted multiplicative ring section.

If the level action is only locally free, let
\[
 \ell=\operatorname{lcm}\{|G_z|:z\in Z\}
\]
for its finite stabilizers. In the compact setting this uses finitely many stabilizer orders. The aid reports a stable complex-cobordism \(MU[1/\ell]\)-splitting and the resulting coefficient-localization statement: for every commutative \(\Z[1/\ell]\)-algebra \(R\), \(\kappa_R^{\mathrm B}\) is additively split surjective. With these coefficients, the finite-isotropy comparison
\[
 q^*:H^*(Z/G;R)\xrightarrow{\ \cong\ }H_G^*(Z;R)
\]
identifies Borel and coarse targets, so \((q^*)^{-1}\kappa_R^{\mathrm B}\) is also split surjective. This is not an assertion of integral coarse-space surjectivity when stabilizers remain nontrivial.

For the integral Borel cokernel
\[
 C=\operatorname{coker}\bigl(\kappa_\Z^{\mathrm B}\bigr),
\]
the supplied localization audit records
\[
 C\otimes_\Z\Z[1/\ell]=0,
 \qquad
 \forall c\in C\ \exists N(c)\ge0:\ \ell^{N(c)}c=0.
\]
Thus \(C\otimes\Z_{(p)}=0\) for every prime \(p\nmid\ell\). Localization alone gives neither \(\ell C=0\) nor one exponent \(N\) valid for every element. In the nonfree integral setting the map \(q^*\) need not be an isomorphism, so restriction by itself does not canonically define a map to integral coarse-space cohomology. The source reports these results as broad progress on the free/finite-stabilizer cases, not a classification of every possible symplectic reduction.
\subsection{Short English statement}
The problem asks when integral Kirwan surjectivity holds. The supplied aid reports a split-surjectivity theorem for free regular levels and a version for finite stabilizers after their orders are inverted. The remaining integral Borel cokernel is elementwise stabilizer-power torsion; this does not supply an integral coarse-space map or a uniform torsion bound.
\subsection{Sources}
\begin{itemize}
\item Local UnsolvedMath registry, permanent ID 20001975, exact source code catalogue entry; canonical \texttt{aim-geometry-notes.json}, \texttt{notes[312]}, tag \texttt{problem}; AIM Workshop Problem Lists, Geometry, Problem 10.1 (R. Goldin, S. Tolman, J. Weitsman). Created and updated 14 August 2026.
\item AIM \emph{Moment maps and surjectivity in various geometries}; \url{https://aimath.org/WWN/momentmaps/momentmaps.pdf}. The supplied verification identifies the original problem; the PDF was not accessed.
\item Pomerleano--Teleman theorem and stable \(MU[1/\ell]\)-splitting, as attributed in the supplied research summary. No title, date, publication identifier, theorem number, or separate URL is supplied, so the exact statement and hypotheses require source verification.
\item Supplied attempt 1: \texttt{partial\_result}, problem status at run \texttt{partially\_solved}, rigor label \texttt{proved}, novelty confidence \texttt{low}; classification \texttt{PARTIAL-PROGRESS}, registry status \texttt{partially\_solved}. Candidate: coefficient-localization and target audit. The summary's ``current theorem'' assessment is dated supplied evidence, not a fresh literature judgment.
\item Ulam AI / UnsolvedMath Contributors; CC BY 4.0. Machine-generated research aid, not peer reviewed; independently verify before citation.
\end{itemize}
\end{document}
